\section{Introduction}
\label{sec:introduction}

Markov constraints on a rooted tree amplify local choices very differently
from constraints on a one-dimensional shift.  A decision made at a vertex is
replicated across exponentially many descendants, while the natural metric
assigns its scale through the size of an entire finite tree.  In this setting,
even one strictly transient graph state can affect a boundary-weighted
dimension: its $d$ children can enter different cyclic phases, and the chosen
phase proportions persist through exponentially large descendant subtrees.
The phenomenon studied here is this finite integer allocation, not arbitrary
reducibility.

Tree-shifts of finite type and Markov hom tree-shifts have an established
entropy theory; see, among others, \citet{aubrunBeal2012},
\citet{petersenSalama2018}, and \citet{petersenSalama2020}.  Reducible
topological entropy need not be the maximum of the component entropies
\citep{banEtAl2021,banEtAl2022}.  More recently,
\citet{banLaiWu2025} developed Hausdorff-dimension theory for irreducible
Markov hom tree-shifts in the metric used below and identified the general
reducible case as a separate problem.  These results set the boundary of the
present article: we neither claim the component-failure mechanism for
topological entropy nor solve the general reducible Hausdorff problem.

Our core is deliberately rigid.  Its alphabet is partitioned into $p$
nonempty phases $V_0,\ldots,V_{p-1}$, and every label in $V_j$ may be followed
by every label in $V_{j+1}$, with indices modulo $p$.  Put
$a_j=|V_j|$ and $c_j=\log a_j$.  For the complete core, direct counting gives
the supporting quantity
\[
 H_j(c)=\frac{d-1}{d^p-1}
 \sum_{t=0}^{p-1}d^{p-1-t}c_{j-t}.
\]
The core dimension is $\min_j H_j(c)$.  We then add one transient state $r$
that sees every core label and has no return edge.  Fixing the phases of its
ordered children gives a composition $m=(m_0,\ldots,m_{p-1})$ of $d$ and the
dimension
\[
 D_1(m;c)=\min_j\frac1d\sum_s m_s H_{s+j}(c).
\]
There are only finitely many phase assignments, so maximizing this expression
over integer compositions gives the dimension of the full one-level shift.
This order of operations matters: the minimum over depth residues remains
inside the integer optimization.

\begin{figure}[t]
  \centering
  \begin{tikzpicture}[
  >=Latex,
  every node/.style={font=\small},
  state/.style={circle,draw=black!70,thick,minimum size=7mm,inner sep=0pt},
  phase/.style={rounded corners,draw=black!60,thick,minimum width=16mm,minimum height=8mm},
  note/.style={align=center,font=\footnotesize}
]
  \node[state,fill=black!8] (r) at (0,0) {$r$};
  \node[note,above=2mm of r] {one transient\\root};

  \node[state,fill=blue!16] (c0) at (-3.0,-1.6) {$V_0$};
  \node[state,fill=orange!22] (c1) at (-1.0,-1.6) {$V_1$};
  \node[state,fill=gray!12] (cmid) at (1.0,-1.6) {$\cdots$};
  \node[state,fill=green!18] (clast) at (3.0,-1.6) {$V_{p-1}$};
  \foreach \child in {c0,c1,cmid,clast} {
    \draw[->,thick] (r) -- (\child);
  }
  \node[note,left=2mm of c0] {$m_0$ children};
  \node[note,right=2mm of clast] {$m_{p-1}$ children};
  \node[note] at (0,-2.35) {$m_s\geq0$, \quad $\sum_s m_s=d$};

  \node[phase,fill=blue!16] (v0) at (-3.0,-3.45) {$V_0$};
  \node[phase,fill=orange!22] (v1) at (-1.0,-3.45) {$V_1$};
  \node[phase,fill=gray!12] (vmid) at (1.0,-3.45) {$\cdots$};
  \node[phase,fill=green!18] (vlast) at (3.0,-3.45) {$V_{p-1}$};
  \draw[->,thick] (v0) to[bend left=12] node[above,font=\scriptsize] {complete} (v1);
  \draw[->,thick,dashed] (v1) -- (vmid);
  \draw[->,thick,dashed] (vmid) -- (vlast);
  \draw[->,thick] (vlast) to[bend left=20] (v0);
  \node[note] at (0,-2.92) {complete cyclic core: $V_j\longrightarrow V_{j+1}$};

  \draw[decorate,decoration={brace,amplitude=5pt},thick]
    (-3.55,-1.15) -- (3.55,-1.15)
    node[midway,above=7pt,note] {ordered children, grouped only for counting};

  \node[draw=black!55,rounded corners,fill=white,align=center,text width=10.8cm]
    at (0,-5.35) {$\displaystyle
      \dim_H X_r
      =\max_{\substack{m_s\in\mathbb Z_{\geq0}\\\sum_s m_s=d}}
       \min_j\frac1d\sum_s m_s H_{s+j}(c)$\\[-1mm]
      \footnotesize The finite feeder chooses the integer phase allocation;
      the infinite core supplies the residue-dependent weights.};
\end{tikzpicture}

  \caption{The unrestricted one-level feeder chooses an integer phase
  allocation $m$ at its frontier.  The complete cyclic core turns each fixed
  allocation into an equiprobable stratum, whose least depth residue controls
  its Hausdorff dimension.  Children are grouped only for counting; the tree
  remains ordered.  The diagram does not represent an arbitrary reducible
  graph.}
  \label{fig:mechanism}
\end{figure}

The complete-core calculation above is an ingredient rather than a novelty
claim about irreducible theory.  The residual contributions are as follows.

\begin{enumerate}[label=(\roman*)]
  \item We derive the exact phase-allocation max--min for the unrestricted
  one-level feeder.  The proof includes both Hausdorff bounds, the finite-union
  step, and the fact that concentrated compositions dominate all core-root
  strata.  The same argument covers only explicitly declared finite sets of
  one-level compositions.

  \item We identify the exact saturation arithmetic.  Every allocation lies
  below the spectral mean $\bar c=p^{-1}\sum_jc_j$, and equality is equivalent
  to constancy of a circular convolution, or equivalently of shifted integer
  products.  Uniform allocation proves $p\mid d$ sufficient.  The converse is
  stated only when every nonzero Fourier mode of $c$ is nonzero; an explicit
  $p=4,d=2$ witness rejects an unconditional converse.

  \item We obtain closed formulas for two phases and extend the allocation to
  the canonical unrestricted $L$-level forced chain.  The exact denominator is
  $d^L$, the optimized dimensions are monotone, and a balanced composition
  gives an explicit $O(d^{-L})$ distance to $\bar c$.  Restricted multilevel
  variants are mentioned only with a stated balanced-access condition.

  \item We exhibit a four-state binary shift for which the sole cyclic
  essential strongly connected component has dimension $\log 2/3$, whereas
  the full reducible shift has dimension $\log 2/2$.  Thus the cyclic-essential-SCC
  maximum fails.  Exact programs replay
  the identities and negative controls, but no finite census substitutes for
  the universal proofs.
\end{enumerate}

The proof strategy is elementary but sensitive to indices and metric scales.
We first identify balls with finite-tree cylinders and prove a matching
cylinder-cover/Frostman lemma.  We then compute periodic weighted limits with
the backward index $j-t$.  Only after the exact one-level optimizer is proved
do we study equality through a circulant operator and discrete Fourier
modes.  Finally, the map $m\mapsto dm$ embeds the depth-$L$ allocation grid in
the depth-$(L+1)$ grid, while balanced integer compositions control the
remaining gap.

The scope is narrow throughout.  We assume positive phase sizes, complete
phase-to-phase blocks, and no return to a feeder.  Arbitrary finite strictly
transient feeder shapes, incomplete blocks, communicating cyclic components,
and nontransient reuse are excluded.  Section~\ref{sec:related} fixes the
source boundary; Sections~\ref{sec:setting}--\ref{sec:deep} prove the main
results; Section~\ref{sec:example} gives the four-state application and exact
verification; and the appendices record endpoints and reproducibility.
